\section{Introduction}\label{sec:intro}
An autonomous agent can continue acting while advice is computed or communicated. When a reply arrives, the agent faces two questions: whether to change its behavior, and whether any resulting improvement can be credited to the information in that reply. The second question is easily obscured by the first. A policy trained with advice features may outperform a policy trained without them because fitting or model selection changes, even when the final deployed decisions do not depend on those features.

Action advising has long distinguished what to teach, when to teach and how a learner uses advice. Teacher quality and student benefit need not have a simple relationship~\cite{fachantidis}; demonstrations, corrections and guidance enter a learner through different mechanisms~\cite{najar}. Confidence-based autonomy regulates requests for demonstrations~\cite{chernova}, while advice reuse and imitation transfer teacher knowledge into the student~\cite{zhu,ilhan}. These results make offline knowledge and online communication distinct parts of an advice system. They also motivate a precise evaluation question: when two advice-labelled pipelines obtain different rewards, does the deployed behavior actually depend on the returned message?

Delayed control sharpens the attribution problem. Action Chunking with Transformers (ACT) predicts sequences from demonstrations~\cite{act}, and real-time chunking addresses asynchronous execution for flow-based policies~\cite{rtc}. Yet a capable complete-episode source need not be a useful occasional advisor to an already competent local controller. At arrival, a selector may change a local action, execute source actions, or reject. A feature-availability ablation can simultaneously change representation, fitted coefficients and validation selection. Its reward difference is therefore a whole-pipeline result unless additional controls establish the role of the runtime message.

We study this issue through a simulation evaluation framework for delayed advice in autonomous decision making. The framework makes the target contrast explicit, retains executable local and null-advice references, records decisions before counterfactual outcomes, and checks native branch replay. Its central empirical case is a positive packet-input-labelled reward contrast whose selected policy is identically zero and always rejects intervention. This is a concrete counterexample to a particular attribution, rather than evidence that advice is generally ineffective. A simple dependency certificate explains why its behavior can be reproduced without a returned packet.

The paper contributes three connected elements. First, it distinguishes fitted-pipeline performance, fixed-policy message dependence and value under a specified channel intervention, with an elementary counterexample and an observed executable null policy. Second, it turns those distinctions into an operational audit: a workflow, an algorithm, pre-outcome decisions, native branch replay and explicit message-generation/fallback rules. Third, it supplies the visual attribution counterexample on two manipulation tasks and direct fixed-policy channel outcomes on the original frozen visual PickPlace system and three state-observation contact tasks, with finite-pool diagnostics that distinguish limited opportunity from harmful selection. Further experiments examine critic data, training variation, delay, sensor noise, request schedules and model-selection tolerance. The contribution is an evaluation framework and its empirical failure mechanisms; standard regression, null models and matched branching are established components, not claimed algorithmic inventions.

The agent-level relevance is methodological. Teacher--student and decision-support architectures are often evaluated through end-to-end reward. Establishing an improvement remains necessary, but identifying its origin requires checking what the deployed system consumes and comparing outcomes under a defined intervention. Our implementation uses a single learner and a fixed computational advisor; it does not demonstrate strategic multi-agent interaction. The abstract distinction applies beyond robotics, while its empirical support here remains limited to simulated manipulation.
